\documentclass[conference]{IEEEtran}

\usepackage{cite}
\usepackage{amsmath,amssymb,amsfonts}
\usepackage{graphicx}
\usepackage[hyphens]{url}
\Urlmuskip=0mu plus 1mu\relax
\usepackage{textcomp}
\usepackage{xcolor}
\usepackage{booktabs}
\usepackage{tagging}

\usepackage{sc26repro}

% The next two lines MUST stay commented for final submission.
%\usetag{explanation}
%\usetag{example}

\begin{document}

\twocolumn[%
{\begin{center}
\Huge
Appendix: Artifact Description/Artifact Evaluation
\end{center}}
]

%%%%%%%%%%%%%%%%%%%%%%%%%%%%%%%%%%%%%%%%%%%%%%%%%%%%%
% AD Appendix
%%%%%%%%%%%%%%%%%%%%%%%%%%%%%%%%%%%%%%%%%%%%%%%%%%%%%

\appendixAD

\section{Overview of Contributions and Artifacts}

\subsection{Paper's Main Contributions}

\begin{description}
\item[$C_1$] We provide a large-scale, trace-grounded
characterization of AI network sensitivity across realistic
workloads, fabrics, and scales up to 4{,}096 GPUs. The study
separates the effects of latency, bandwidth, and jitter and
identifies the operating regimes in which each limits performance.
\item[$C_2$] We develop and validate a compositional
linear-programming method with signature-based caching. The method
preserves workload dependencies while making production-scale
sensitivity studies practical.
\item[$C_3$] We derive latency-bandwidth performance envelopes and
use them to identify safe operating regions, performance cliffs,
and useful network provisioning points.
\end{description}

\subsection{Computational Artifacts}

The paper is accompanied by two artifacts. Artifact $A_1$ contains
the analysis code, compact result data, reproduction drivers,
plotting scripts, manifests, and the final plots. Artifact $A_2$ is
the companion collection of execution traces used as inputs to the
analysis.

\begin{description}
\item[$A_1$] \url{https://github.com/tommasobo/SC_Tracing}
  (branch \texttt{artifact\_freeze})
\item[$A_2$] \url{http://storage2.spcl.ethz.ch/traces/ai/}
\end{description}

The authors intend to request the Artifacts Available, Artifacts
Evaluated-Functional, and Results Reproduced badges. Persistent
DOIs will be assigned at the SC26 artifact-freeze stage. No DOI has
been assigned at the time of this submission.

\begin{center}
\small
\begin{tabular}{rll}
\toprule
Artifact & Contributions & Related paper elements \\
\midrule
$A_1$ & $C_1,C_2,C_3$ & Figs.\ 1, 3--10 \\
$A_2$ & $C_1,C_2$ & Table II, Figs.\ 3--6 \\
\bottomrule
\end{tabular}
\end{center}

\section{Artifact Identification}

\newartifact

\artrel

Artifact $A_1$ contains the following components:

\begin{itemize}
\item \texttt{solver/}: Composite-LP and Monolithic-LP analysis.
\item \texttt{tools/}: the NCCL-to-GOAL conversion code and the
  LogGOPSim source.
\item \texttt{pipeline/}: drivers for trace conversion, LP analysis,
  LogGOPSim replay, and selected sensitivity sweeps.
\item \texttt{data/} and \texttt{results/}: compact inputs and
  outputs for the figure workflows.
\item \texttt{scripts/}: one plotting workflow for each
  reproducible paper figure.
\item \texttt{local\_artifact/}: provenance manifests and checks for
  the recovered Figure 5 and Figure 6 inputs.
\item \texttt{reproduce\_quick.sh} and
  \texttt{reproduce\_full.sh}: the reviewer entry points.
\end{itemize}

The trace-to-analysis path supports $C_1$ and $C_2$. Its outputs
are the sensitivity envelopes, validation comparisons, memory
measurements, and jitter results that support $C_3$. Figure 2 is a
method diagram and has no computational output.

\artexp

The quick workflow redraws Figures 1 and 3--10 from
the compact data. It should complete in a few minutes on a laptop
and produce PDF and PNG files under \texttt{figures/}. The layout,
curves, relative ordering, and main breakpoints should agree
closely with the paper.

The full workflow additionally exercises affordable analysis and
trace-processing paths. It supports local execution by default and
Slurm execution on Alps or a similar cluster. Long or high-memory
stages are disabled unless the user supplies
\texttt{--expensive\_run}. In particular, the Grok 4{,}096-GPU
experiment is not run by default.

The expected scientific results are:

\begin{itemize}
\item Composite-LP and the validation baselines have close runtime
  predictions for the configurations where the baselines are
  feasible.
\item The workload curves have different sensitivity regimes.
  Llama training is bandwidth-sensitive over a substantial range,
  Grok at large scale has a broad compute-dominated region, and
  vLLM inference has a broad memory-dominated region.
\item The full-trace Monolithic-LP memory requirement grows quickly
  with scale, while the compositional method remains practical.
\item Sparse jitter has little effect on the studied Grok workload,
  while perturbing a large fraction of messages produces a clear
  slowdown.
\end{itemize}

The repository includes a reproduction report with quantitative
comparisons and known differences. The most important limitations
are also stated here. The released 128-GPU Llama source is named
\texttt{Llama7B\_N32\_GPU128}, although the paper labels the panel
as Llama 70B. The matching scale and the close regenerated curve
indicate a paper-label error. The closest recovered Figure 5
Composite curve is raw-derived and differs from the paper CSV by
at most 0.014853\%. Cold runs from the recovered Llama N4/GPU16
traces do not match it as closely. The recovered Figure 6 vLLM raw
trace follows the complete conversion and replay path, but a
paper-era token-window reduction step is not available. The
default vLLM plot therefore uses the compact paper data, and the
report records the raw-replay difference.

\arttime

Times below are estimates for an ordinary workstation unless a
cluster is stated.

\begin{itemize}
\item Artifact setup: 5--10 min for cloning, creating a Python
  environment, and installing the plotting dependencies.
\item Quick execution: less than 5 min for
  \texttt{./reproduce\_quick.sh}.
\item Quick analysis: 10--20 min to run the checks and compare the
  generated plots with the paper and reference files.
\item Full execution without expensive stages: up to 6 h. Each
  selected figure workflow is designed to finish within about one
  hour. Independent tasks may be submitted concurrently on Slurm.
\item Full analysis: about 30 min for manifest checks, plot
  inspection, and review of the generated report.
\end{itemize}

The quick path is sufficient to exercise the packaged artifact.
The non-expensive full path fits within the SC26 eight-hour review
budget. Grok at 4{,}096 GPUs and any stage expected to exceed
45 min or require substantial memory remain opt-in and are not part
of the recommended reviewer run.

\artin

\artinpart{Hardware}

The quick workflow requires a laptop or workstation with about
1\,GiB of free RAM and 1\,GiB of free disk space. It does not
require a GPU.

The full local workflow should use scratch storage and a host with
at least 16\,GiB of RAM. Raw NSYS conversion and larger GOAL files
may require more memory and disk. The Slurm path was prepared for
the CSCS Alps system, whose compute nodes use NVIDIA GH200
superchips and the Slingshot-11 interconnect. Large trace runs must
use compute nodes rather than login nodes.

The opt-in Grok 4{,}096-GPU path and large Monolithic-LP cases can
require hundreds of GiB to multiple TiB of memory or very large
temporary files. Reviewers should not enable this path without an
appropriate allocation.

\artinpart{Software}

The default artifact requires:

\begin{itemize}
\item Python 3.10 or later, \url{https://www.python.org/}.
\item NumPy 1.21 or later,
  \url{https://numpy.org/}.
\item pandas 1.3 or later,
  \url{https://pandas.pydata.org/}.
\item Matplotlib 3.5 or later,
  \url{https://matplotlib.org/}.
\item Bash 4.4 or later,
  \url{https://www.gnu.org/software/bash/}.
\end{itemize}

The optional LogGOPSim build requires a C++ compiler,
\texttt{gengetopt}, and \texttt{re2c}:
\url{https://gcc.gnu.org/},
\url{https://www.gnu.org/software/gengetopt/}, and
\url{https://re2c.org/}.

LP regeneration requires Gurobi 10.0 or later,
\url{https://www.gurobi.com/}. A free academic license is
sufficient. Slurm execution requires an installed Slurm client,
\url{https://slurm.schedmd.com/}. Conversion from raw
\texttt{.nsys-rep} inputs requires NVIDIA Nsight Systems with
SQLite export support,
\url{https://developer.nvidia.com/nsight-systems}.

\artinpart{Datasets / Inputs}

The compact CSV and JSON inputs required for the default workflows
are stored in the repository. They include sensitivity sweeps,
hardware reference points, memory results, jitter results, and
provenance metadata. These files are small enough for a normal Git
checkout.

Artifact $A_2$ provides larger GOAL, BIN, and NSYS inputs. A
reviewer does not need the complete collection. Selected trace
directories can be downloaded separately. Raw and intermediate
files can be several GiB, while the largest Grok case produces
temporary data on the TiB scale.

Figures 1, 7, 8, 9, and 10 use existing compact data by design.
The default Figure 5 plot uses the closest recovered raw-derived
Composite curve together with existing Monolithic and LGS data.
The Grok and vLLM portions of Figure 6 use compact paper data
because of the limitations described under Expected Results. All
plotting scripts are included.

\artinpart{Installation and Deployment}

\begingroup\small
\begin{verbatim}
git clone -b artifact_freeze \
  https://github.com/tommasobo/SC_Tracing.git
cd SC_Tracing
python3 -m venv .venv
source .venv/bin/activate
python3 -m pip install -r requirements.txt
\end{verbatim}
\endgroup

No compilation is needed for the quick workflow. The optional
LogGOPSim build is performed by
\texttt{pipeline/build\_tools.sh}. On a Slurm cluster, clone or copy
the repository to scratch and keep raw and temporary files there.
The full workflow also requires:

\begingroup\small
\begin{verbatim}
python3 -m pip install -r requirements-dev.txt
\end{verbatim}
\endgroup

This includes the solver and test dependencies. A configured
Gurobi license is required for fresh LP sweeps.

\artcomp

The workflow consists of five tasks:

\begin{itemize}
\item $T_1$: validate compact inputs and provenance manifests.
\item $T_2$: choose the quick plot-only path or the full path.
\item $T_3$: for enabled full-path experiments, convert or load
  traces, build dependency metadata, and run Composite-LP,
  Monolithic-LP, or LogGOPSim as required.
\item $T_4$: transform result tables and generate the paper plots.
\item $T_5$: check the expected output set and write the
  reproduction report.
\end{itemize}

The quick dependency chain is
$T_1 \rightarrow T_2 \rightarrow T_4 \rightarrow T_5$.
The full dependency chain is
$T_1 \rightarrow T_2 \rightarrow T_3 \rightarrow T_4
\rightarrow T_5$. Independent Figure 3--5 tasks may run in
parallel. The two Figure 6 solver sweeps are serialized by the
Slurm launcher to stay within the available concurrent Gurobi
license slots.

The compact-data path is deterministic and needs one run. Hardware
reference data and stochastic experiments retain the repetitions
used for the paper in the packaged tables. Fresh stochastic or
hardware runs must retain their seed and run count in their output
metadata.

The two top-level commands are:

\begingroup\small
\begin{verbatim}
./reproduce_quick.sh
./reproduce_full.sh
\end{verbatim}
\endgroup

The full script runs locally by default. Useful options are
\texttt{--scratch DIR}, \texttt{--workers N}, \texttt{--slurm},
and \texttt{--dry-run}. The \texttt{--expensive\_run} flag is the
only way to enable Grok 4k and other long or high-memory stages.

\artout

Successful execution produces plots under \texttt{figures/} and
processed tables, logs, and execution manifests under the selected
scratch directory. The final reproduction report is under
\texttt{docs/}. The manifest and artifact checks must pass. The
evaluator should compare the generated PDFs with the paper or the
supplied side-by-side reference.

The comparison should focus on curve shape, method ordering,
breakpoints, and the conclusions listed under Expected Results.
Exact pixel equality is not required. Any remaining numeric
difference, including the known Figure 5 and vLLM Figure 6 cases,
must be recorded in the final report rather than hidden by
substituting another workload.

\newartifact

\artrel

Artifact $A_2$ contains the trace inputs used by the analysis. It
supports $C_1$ by making the communication workloads inspectable
and reusable. It supports $C_2$ by providing GOAL schedules that
can be replayed by LogGOPSim and analyzed by the code in $A_1$.

\artexp

A downloaded GOAL trace should parse and replay successfully with
the tools in $A_1$. A selected moderate-size trace should produce
a runtime point or sensitivity curve with the same qualitative
behavior as the corresponding compact result.

The directory name is part of the trace provenance. In particular,
the 128-GPU Llama directory identifies a 7B model. It must not be
renamed to imply that it is a 70B trace. The exact vLLM 8B paper
input is represented by the recovered local provenance manifest
and compact results in $A_1$; the currently public vLLM 70B input
is a different workload and must not be used as a substitute.

\arttime

For one selected moderate trace:

\begin{itemize}
\item Artifact setup: 5--30 min to download the selected files,
  depending on network bandwidth.
\item Artifact execution: about 5--45 min for conversion or replay.
\item Artifact analysis: about 10 min to compare the generated
  runtime points with the compact reference.
\end{itemize}

Downloading the complete trace collection or processing Grok 4k
is outside the recommended reviewer path.

\artin

\artinpart{Hardware}

Downloading and inspecting one selected trace requires a
workstation with sufficient scratch space. Moderate GOAL replays
need several GiB of RAM. Large Grok traces require cluster-scale
scratch and memory and are not enabled by default.

\artinpart{Software}

The downloader requires \texttt{wget} or \texttt{curl}. GOAL replay
uses the LogGOPSim source included in $A_1$. Raw NSYS conversion
uses NVIDIA Nsight Systems and the NCCL converter in $A_1$.

\artinpart{Datasets / Inputs}

$A_2$ is the input dataset. The relevant public top-level
directories are \texttt{llama/}, \texttt{grok/}, and
\texttt{vllm/}. Each directory includes model and scale information
in its name. The manifests in $A_1$ record the input identity and
hashes used during the final reproduction work.

\artinpart{Installation and Deployment}

Download only the trace needed for the selected evaluation. For
example, the public 16-GPU Grok GOAL trace is under:

\begingroup\small
\begin{verbatim}
http://storage2.spcl.ethz.ch/traces/ai/grok/
  Grok314B_N4_GPU16_TP4_PP1_CP1_VP1_
  EP4_ETP4_GBS256/grok.goal
\end{verbatim}
\endgroup

Place downloaded files in a scratch directory. Do not copy large
raw traces into the Git checkout.

\artcomp

The data workflow is:
\begin{center}
\small
NSYS $\rightarrow$ SQLite $\rightarrow$ GOAL/BIN\\
$\rightarrow$ LP or LogGOPSim $\rightarrow$ CSV
$\rightarrow$ plot.
\end{center}
Prebuilt GOAL files enter at the third stage. The raw-input and
GOAL workflows, their metadata requirements, and the validated
exceptions are documented in $A_1$.

\artout

The expected output is a GOAL replay runtime or sensitivity CSV.
The evaluator should keep the exact source URL, file size, and hash
with the output. The generated result can then be passed to the
plotting stage in $A_1$.

%%%%%%%%%%%%%%%%%%%%%%%%%%%%%%%%%%%%%%%%%%%%%%%%%%%%%
% AE Appendix
%%%%%%%%%%%%%%%%%%%%%%%%%%%%%%%%%%%%%%%%%%%%%%%%%%%%%

\newpage
\appendixAE

\arteval{1}

\artin

Clone the \texttt{artifact\_freeze} branch and install the Python
requirements using the commands in the AD. From the repository
root, first inspect the available work:

\begingroup\small
\begin{verbatim}
./reproduce_full.sh --dry-run
\end{verbatim}
\endgroup

The dry run lists the enabled tasks without starting experiments.
Confirm that \texttt{data/}, \texttt{scripts/},
\texttt{reproduce\_quick.sh}, and \texttt{reproduce\_full.sh} are
present. Use scratch storage for all nontrivial runs.

For a local run, choose a scratch directory with:

\begingroup\small
\begin{verbatim}
mkdir -p /path/to/scratch/sc26-provisioning
\end{verbatim}
\endgroup

On Alps, create the scratch directory first and use the Slurm mode
described below. Do not launch trace conversion, LP sweeps, or
LogGOPSim replays on a login node.

\artcomp

\emph{Quick evaluation.} Run:

\begingroup\small
\begin{verbatim}
./reproduce_quick.sh
\end{verbatim}
\endgroup

This validates the compact inputs and redraws Figures 1 and 3--10.
The tasks are: load the committed CSV and JSON inputs,
check their schemas, call the relevant plotting scripts, and check
that the expected PDF and PNG outputs were created. No solver,
large trace, GPU, or cluster allocation is used.

\emph{Full local evaluation without expensive stages.} Run:

\begingroup\small
\begin{verbatim}
./reproduce_full.sh \
  --scratch /path/to/scratch/sc26-provisioning \
  --workers 4
\end{verbatim}
\endgroup

The local driver runs the bounded tasks one after another. Solver
workers within a task use the requested worker count. The core task
redraws the figures and runs the manifest, artifact, and test
checks. Each subsequent task writes result tables, comparison
metrics, and a JSON execution manifest. It does not run Grok 4k or
any stage marked as long or high-memory.

\emph{Full Slurm evaluation.} On Alps or a similar cluster, run:

\begingroup\small
\begin{verbatim}
./reproduce_full.sh --slurm \
  --scratch /path/to/scratch/sc26-provisioning
\end{verbatim}
\endgroup

The driver submits experiments as separate batch jobs. Independent
jobs run concurrently, while the two Figure 6 solver sweeps are
serialized to respect Gurobi license limits. Each job writes its
own log and JSON execution manifest. Monitor the submitted job IDs
and inspect every manifest after completion. The core job performs
the compact-data redraw. Inspect the generated job specifications
first with \texttt{--dry-run}.

\emph{Opt-in expensive evaluation.} The command below is documented
for completeness and is not part of the recommended review:

\begingroup\small
\begin{verbatim}
./reproduce_full.sh --slurm --expensive_run \
  --grok-analysis-dir /scratch/grok/N1024/analysis \
  --scratch /large/scratch/sc26-provisioning
\end{verbatim}
\endgroup

Use it only after confirming the time, memory, storage, and compute
allocation. The flag enables Grok 4k and other stages that may
exceed 45 min or require substantial memory.

\artout

After the quick run, confirm that plots for Figures 1 and 3--10
exist under \texttt{figures/}. After a full run, inspect
the task log and the final report under \texttt{docs/}. Run the
artifact checker documented by the top-level README and verify that
it reports no missing required file or malformed compact dataset.

The expected comparisons are:

\begin{itemize}
\item Figures 1, 7, 8, 9, and 10 are redraws from existing data. Their
  plots should agree with the paper within plotting precision.
\item Figures 3 and 4 should preserve the paper's method ordering,
  sensitivity trends, and main breakpoints. Small configuration
  differences are listed in the report.
\item Figure 5 uses the closest recovered raw-derived Composite
  result, whose maximum relative difference from the paper CSV is
  0.014853\%. The Monolithic and LGS series use existing compact
  data. Cold recovered-trace alternatives and their larger
  differences are documented in the report.
\item Figure 6 should reproduce the plotted paper curves from
  compact data. Fresh Llama metadata analysis closely matches the
  packaged curve, subject to the 7B versus 70B label issue. The
  Grok 4k and vLLM paper curves use compact data by default.
\end{itemize}

For the Results Reproduced assessment, compare behavior rather
than PDF pixels. Check the curve shapes, the ordering of the
methods, the location of the major sensitivity transitions, and
the provisioning conclusions. Use the supplied paper-versus-result
comparison PDF and metrics tables for quantitative checks. The
known differences must remain visible in the report.

\arteval{2}

\artin

Select one moderate trace rather than the full archive. The
16-GPU Grok GOAL trace is suitable for a functional check:

\begingroup\small
\begin{verbatim}
mkdir -p /path/to/scratch/sc26-trace
cd /path/to/scratch/sc26-trace
wget http://storage2.spcl.ethz.ch/traces/ai/grok/\
Grok314B_N4_GPU16_TP4_PP1_CP1_VP1_\
EP4_ETP4_GBS256/grok.goal
\end{verbatim}
\endgroup

Return to the $A_1$ checkout and build LogGOPSim:

\begingroup\small
\begin{verbatim}
pipeline/build_tools.sh
\end{verbatim}
\endgroup

\artcomp

Replay the selected GOAL file at a representative latency:

\begingroup\small
\begin{verbatim}
python3 pipeline/run_lgs.py \
  --goal /path/to/scratch/sc26-trace/grok.goal \
  --L 4000
\end{verbatim}
\endgroup

This path parses the GOAL schedule, converts it to the simulator's
binary format if needed, and runs LogGOPSim. The output reports the
predicted workload runtime. The full driver uses the same stages
for its supported trace cases and stores reusable conversion files
under the selected scratch directory.

Do not substitute the public vLLM 70B trace for the paper's vLLM
8B trace. Do not run the 4{,}096-GPU Grok trace during a routine
evaluation.

\artout

Record the source URL, downloaded file size, hash, LogGOPSim
parameters, and reported runtime. Compare the result with the
matching compact reference or the tolerance stated in the final
reproduction report. A successful parse and replay exercises the
published trace, the simulator build, and the trace-to-runtime
portion of the method. Use $A_1$ for the final plotting and
paper-level comparison.

\end{document}
